\documentclass[12pt]{article}
\usepackage[margin=0.72in]{geometry}
\usepackage{lmodern}
\usepackage{microtype}
\usepackage{multicol}
\usepackage{fancyhdr}
\usepackage{titlesec}
\usepackage{enumitem}
\usepackage{xcolor}
\usepackage[hidelinks]{hyperref}

\setlength{\columnsep}{0.30in}
\setlength{\parindent}{1.05em}
\setlength{\parskip}{0.24em}
\setlength{\headheight}{15pt}
\titleformat{\section}{\large\bfseries\scshape}{}{0pt}{}
\titleformat{\subsection}{\normalsize\bfseries}{}{0pt}{}
\pagestyle{fancy}
\fancyhf{}
\fancyhead[L]{Bacon's Four Idols}
\fancyhead[R]{Reading Handout}
\fancyfoot[C]{\thepage}
\renewcommand{\headrulewidth}{0.4pt}

\newcommand{\focusbox}[1]{%
\vspace{0.4em}\noindent\colorbox{gray!12}{\begin{minipage}{0.96\linewidth}#1\end{minipage}}\vspace{0.5em}}

\begin{document}
\begin{center}
{\LARGE\bfseries Do We Love Shadows?}\par\vspace{0.25em}
{\large\scshape Reading 3: Francis Bacon and the Four Idols}\par\vspace{0.45em}
\rule{0.82\textwidth}{0.4pt}\par\vspace{0.7em}
\end{center}

\section*{Vocabulary Preview}
\begin{multicols}{2}
\begin{itemize}[leftmargin=*, itemsep=0.18em]
\item \textbf{Idol}: For Bacon, a false image or mistaken idea lodged in the mind.
\item \textbf{Understanding}: The mind's power to think, judge, and make sense of things.
\item \textbf{Preoccupied}: Already filled or taken over before something new can enter.
\item \textbf{Notion}: An idea, concept, or mental picture.
\item \textbf{Axiom}: A general principle used as a starting point for reasoning.
\item \textbf{Induction}: Reasoning from particular observations toward a general conclusion.
\item \textbf{Tribe}: Here, the whole human race.
\item \textbf{Den}: Bacon's word for each person's private cave of habits and experience.
\item \textbf{Market}: Public life, conversation, argument, trade, and shared language.
\item \textbf{Dogma}: A teaching or system of belief treated as settled truth.
\item \textbf{Demonstration}: A proof or method of showing that something is true.
\item \textbf{Inveterate}: Deeply settled by long habit or tradition.
\end{itemize}
\end{multicols}

\section*{Introduction}
Francis Bacon was an English philosopher and statesman who lived from 1561 to 1626. He wrote during a time when many thinkers were questioning old authorities and seeking a better method for discovering truth about nature. In \textit{Novum Organum}, Bacon argues that the search for truth is difficult not only because the world is complicated, but because the human mind itself is easily distorted.

Plato gives us the image of prisoners mistaking shadows for reality. Homer then shows another danger: a shadow may be attractive because it sounds beautiful and promises wisdom. Bacon asks a related question: \textbf{what makes the mind produce, accept, and defend shadows in the first place?} He calls these false images \textit{idols}. Before people can reason well, Bacon says, they must learn to recognize the idols that already possess the mind.

\focusbox{\textbf{Source note:} Adapted and modernized from Francis Bacon, \textit{Novum Organum}, Book I, Aphorisms 38--44, in Joseph Devey's public-domain edition. This version preserves Bacon's sequence, aphoristic structure, four categories, and central images while using modern student-facing English.}

\begin{multicols}{2}
\section*{Reading}

\subsection*{38. Idols Already in the Mind}
The idols and false notions that have already taken possession of the human understanding are deeply rooted there. They do not merely fill people's minds so that truth can hardly enter. Even when truth does gain entrance, these idols return and trouble us again, especially when we try to begin the sciences anew.

Because of this, human beings must be warned ahead of time. Once they know the danger, they must guard themselves with all possible care against these idols and false notions.

\subsection*{39. Four Kinds of Idols}
Four kinds of idols beset the human mind. For the sake of distinction, I give them four names.

The first are \textbf{Idols of the Tribe}. The second are \textbf{Idols of the Den}. The third are \textbf{Idols of the Market}. The fourth are \textbf{Idols of the Theatre}.

\subsection*{40. The Remedy}
The only proper remedy for these idols is the formation of true notions and axioms by true induction. In other words, the mind must learn to build its ideas carefully from reality, not merely from habit, opinion, or inherited words.

Still, it is very useful simply to point the idols out. The doctrine of idols helps us interpret nature in much the same way that the study of false arguments helps ordinary logic. If we know the traps, we are less likely to fall into them.

\subsection*{41. Idols of the Tribe}
The \textbf{Idols of the Tribe} are rooted in human nature itself. They belong to the very tribe or race of man.

People falsely say that human sense is the measure of things. In truth, the perceptions of both the senses and the mind refer more to human beings than to the universe itself. The human mind is like an uneven mirror. It receives rays from things, but it mixes its own nature with their nature. As a result, it distorts and disfigures what it sees.

These idols are not merely the private mistakes of one foolish person. They arise from the shared weaknesses of human beings as human beings.

\subsection*{42. Idols of the Den}
The \textbf{Idols of the Den} belong to each individual person. Besides the errors common to the human race, each person has his own private den or cave.

This private cave intercepts and corrupts the light of nature. It may be shaped by a person's particular disposition, by education, by conversation with others, by reading, by the authority of people he admires, or by the different impressions made on a mind that is already occupied and predisposed. One mind may be restless and biased; another may be calmer and more even. In each case, the human spirit is variable and confused, almost as if moved by chance.

Heraclitus said well that human beings search for knowledge in smaller private worlds, not in the greater world they share in common.

\subsection*{43. Idols of the Market}
There are also idols formed by the mutual dealings and society of human beings. These I call \textbf{Idols of the Market}, because they arise from the commerce, association, and conversation of people with one another.

Human beings converse by means of language. But words are often formed according to common opinion, not according to the true nature of things. Because words are badly and carelessly formed, they create a remarkable obstruction to the mind.

Definitions and explanations can sometimes help learned people protect themselves, but they cannot completely cure the problem. Words still plainly force the understanding, throw everything into confusion, and lead people into empty controversies and countless errors.

\subsection*{44. Idols of the Theatre}
Last, there are idols that have crept into people's minds from the various dogmas of philosophical systems and from faulty rules of demonstration. These I call \textbf{Idols of the Theatre}.

I use that name because many accepted or imagined systems of philosophy are like plays brought onto a stage and performed. They create fictional and theatrical worlds. They may be impressive, orderly, and persuasive, but they are still made worlds rather than reality itself.

I do not mean only the systems that already exist, or only the philosophies and sects of the ancient world. Many more such plays can still be composed, and they can be made to agree with one another even while they lead people into opposite errors. Nor do I mean only whole systems. Many smaller elements and axioms of the sciences have also become deeply rooted through tradition, easy belief, and neglect.

Therefore, each kind of idol must be discussed more fully and distinctly, so that the human understanding can be guarded against them.

\section*{Reading Focus}
\begin{enumerate}[leftmargin=*, itemsep=0.2em]
\item Underline one sentence that explains why truth has trouble entering the mind.
\item Put a \textbf{T} in the margin beside Bacon's definition of Idols of the Tribe.
\item Put a \textbf{D} in the margin beside Bacon's definition of Idols of the Den.
\item Put an \textbf{M} in the margin beside Bacon's explanation of how words confuse thought.
\item Put a \textbf{Th} in the margin beside Bacon's reason for calling some idols \textit{Idols of the Theatre}.
\item Box one sentence that best connects Bacon back to Plato's cave or Homer's Sirens.
\end{enumerate}

\section*{Four Idols Card Sort}
For each situation below, identify which idol is most likely at work. Then write one sentence explaining your choice.

\begin{enumerate}[leftmargin=*, itemsep=0.35em]
\item A student believes a rumor because many people repeat the same catchy phrase.\newline \textbf{Idol:} \underline{\hspace{0.7in}}\newline \textbf{Reason:} \underline{\hspace{1.9in}}
\item A person rejects a new idea immediately because it does not fit what his favorite author taught him.\newline \textbf{Idol:} \underline{\hspace{0.7in}}\newline \textbf{Reason:} \underline{\hspace{1.9in}}
\item A person sees three examples of something and assumes it must always happen that way.\newline \textbf{Idol:} \underline{\hspace{0.7in}}\newline \textbf{Reason:} \underline{\hspace{1.9in}}
\item Two students hear the same speech. One notices every insult; the other notices every joke, because each is used to paying attention to different things.\newline \textbf{Idol:} \underline{\hspace{0.7in}}\newline \textbf{Reason:} \underline{\hspace{1.9in}}
\end{enumerate}

\section*{Window Note}
Create a four-pane window note titled \textbf{``Where Do Shadows Come From?''}

\begin{itemize}[leftmargin=*, itemsep=0.2em]
\item \textbf{Pane 1: Tribe} --- Draw a shadow that comes from ordinary human weakness. Add one sentence explaining the mistake.
\item \textbf{Pane 2: Den} --- Draw a shadow that comes from someone's private background, training, or habits. Add one sentence explaining the mistake.
\item \textbf{Pane 3: Market} --- Draw a shadow that comes from confusing or powerful words. Add one sentence explaining the mistake.
\item \textbf{Pane 4: Theatre} --- Draw a shadow that comes from a whole system, authority, or ideology. Add one sentence explaining the mistake.
\end{itemize}

\section*{Connection to Plato and Homer}
Answer in one clear paragraph: \textbf{How does Bacon help explain either Plato's cave or Homer's Sirens?} Your paragraph should state a claim, name at least one of Bacon's idols, explain how that idol could keep someone trapped in shadows or drawn toward a dangerous song, and connect your answer to education.

\section*{Handout Summary}
Finish by writing a short summary of this handout. In three to five sentences, explain Bacon's main warning about the human mind and name the four idols he says we must learn to recognize. End with one sentence explaining why Bacon belongs after Plato and Homer.

\end{multicols}
\end{document}
